\documentclass{beamer}
\usetheme{Boadilla}
\usepackage{graphicx}
\usepackage[utf8]{inputenc}        
\usepackage[T2A]{fontenc}          
\usepackage[english,russian]{babel}


\title[GaussQE]{Оптимизация ChrobELIAS с помощью GaussQE}
\author[Самодуров Н., Стариков Е.]{Самодуров Никита, Стариков Егор}
\institute[]{
  Научный руководитель: Белянин Георгий \\
  СПбГУ
}
\date{29.07.2026}

\begin{document}

\maketitle

\begin{frame}{Введение}
    \begin{itemize}
      \item ChrobELIAS — это специализированный SMT-решатель (Satisfiability Modulo Theories), предназначенный для проверки выполнимости логических формул в рамках экспоненциально-линейной целочисленной арифметики с регулярными ограничениями. 
      \item GaussQE — алгоритм элиминации кванторов существования.
    \end{itemize}
\end{frame}

\begin{frame}{Акутальность}
    \begin{itemize}
        \item SMT-решатели применяются в таких сферах, как: 
        \begin{itemize}
            \item Верификация ПО и аппаратного обеспечения
            \item Автоматическое тестирование и поиск уязвимостей
            \item Решение задач планировани
        \end{itemize}
        \item Уникальность ChrobELIAS:
        \begin{itemize}
            \item Поддерживает экспоненциально-линейную целочисленную арифметику
            \item Поддерживает регулярные ограничения на строки
        \end{itemize}
    \end{itemize}
\end{frame}

\begin{frame}{Цель и задачи}
    \begin{itemize}
      \item \textbf{Цель:} оптимизация ChrobELIAS путем дополнительной обработки формул с помощью GaussQE.
      \item \textbf{Задачи:}
      \begin{itemize}
        \item Реализация алгоритма GaussQE
        \item Добавление симплификации в ChrobELIAS
        \item Сравнительный анализ производительности
      \end{itemize}
    \end{itemize}
\end{frame}

\begin{frame}{GaussQE}
    \begin{itemize}
        \item Метод элиминации Гаусса-Жордана для целочисленного программирования — это алгоритм, принимающий систему линейных уравнений/неравенств вида $\exists x \, \phi$, который с помощью введения slack-переменных и недетерменированных подстановок переменных избавляется от переменных, стоящих под квантором существования.
    \end{itemize}
\end{frame}

\begin{frame}{Сравнение}
    
\end{frame}

\begin{frame}{Список литературы}
    \begin{thebibliography}{99}
    
    \bibitem{chistikov2024}
    D. Chistikov, A. Mansutti, M.~R. Starchak.
    \newblock Integer Linear-Exponential Programming in NP by Quantifier Elimination.
    \newblock \emph{arXiv:2407.07083}, 2024.
    \end{thebibliography}
\end{frame}

\end{document}
